
\chapter{Additional Experimental Details}\label{app:details}

\section{XGBoost Hyperparameters}

\Cref{tab:xgb-params} lists the XGBoost hyperparameters used across all experiments.
The key constraint is \texttt{max\_depth=2}, which ensures that each tree uses at most two features and can be exactly decomposed into main effects and pairwise interactions.

\begin{table}[H]
\centering
\small
\caption{XGBoost hyperparameters used in all experiments.}
\begin{tabular}{lcc}
\toprule
\textbf{Parameter} & \textbf{German Credit} & \textbf{Stress-Test} \\
\midrule
\texttt{n\_estimators} & 120 & 60 \\
\texttt{max\_depth} & 2 & 2 \\
\texttt{learning\_rate} & 0.1 & 0.1 \\
\texttt{reg\_lambda} (L2) & 0 & 0 \\
\texttt{subsample} & 1.0 & 1.0 \\
\texttt{colsample\_bytree} & 1.0 & 1.0 \\
\bottomrule
\end{tabular}
\label{tab:xgb-params}
\end{table}

\section{Purification Convergence}

In all experiments, the purification algorithm converges to the specified tolerance ($\varepsilon = 10^{-10}$) within 5--15 sweeps per interaction matrix.
The majority of mass transfer occurs in the first sweep, consistent with the theoretical convergence guarantees described in \cref{sec:mass-moving}.
No convergence failures were observed across any experimental run.

\section{Reproducibility}

All experiments are fully reproducible.
The complete source code, experiment scripts, and result generation pipeline are available at:
\begin{center}
\url{https://github.com/HaykTarkhanyan/fanova_purification}
\end{center}

To reproduce all experiments and generate result files:
\begin{verbatim}
pip install -r requirements.txt
python run_all_experiments.py
\end{verbatim}

The stress-test experiments use 3 random seeds per $(\rhoVal, \text{function})$ combination, and all results are reported as means over these seeds.

\chapter{Use of AI Tools}\label{app:ai-tools}

In the preparation of this report and the accompanying implementation, AI-assisted tools were used for the following purposes:
\begin{itemize}
    \item \textbf{Code development:} AI coding assistants were used to help write and debug Python code, particularly for the XGBoost tree decomposition and plotting utilities.
    \item \textbf{Writing assistance:} AI tools were used to assist with drafting and editing the text of this report, including structuring arguments, improving clarity, and ensuring consistent notation.
    \item \textbf{Literature search:} AI tools were used to identify relevant prior work and verify bibliographic details.
\end{itemize}

All AI-generated content was reviewed, verified, and edited by the author.
The scientific contributions---the experimental design, the choice of metrics, the interpretation of results, and the critical analysis of ethical concerns---reflect the author's own judgment.
